% Analytic coverage limits; numerical values are generated by
% verify_coverage_limits.py and written to results/coverage_limits.json.
\subsection{What is, and is not, proved by the coverage limits}
\label{sec:coverage-limits}
Let $R=1800$~m and let $q$ denote the receive radius.  A detector at
radius $d$ subtends on $|x|=R$ an angular half-width $\alpha$ satisfying
\[
 \cos\alpha={R^2+d^2-q^2\over 2Rd}.
\]
Maximising over $d$ gives $\alpha\leq\arcsin(q/R)$.  Consequently, at the
conservative $q=900$~m, six detectors have total boundary-arc capacity at
most $2\pi$.  Equality requires every detector to attain the same extremum
and the six arcs to tile without overlap; the centres then lie on the ring
$d=900\sqrt3$, whose distance from the origin exceeds $900$.  Thus the
origin is uncovered in the equality case, and seven points are necessary for
this $q=900$ certificate.  The origin plus the six-ring construction proves
sufficiency.

At the minimum guaranteed receive radius $q=1000$~m, the same
\emph{fixed} seven-point set is irreducible: deleting a ring point leaves five
ring points, with boundary capacity at most $10\arcsin(5/9)<2\pi$; deleting the
origin leaves it uncovered because $900\sqrt3>1000$.  This is a
fixed-subset statement, not a claim that every arbitrary six-point placement
fails at $q=1000$.

For Q4, put $s=990$~m and use the actual 31 points returned by
\texttt{directional\_waypoints}.  The script calls this function, reconstructs
integer axial coordinates, and enumerates every elementary triangle adjacent
to every one of the 31 points.  For each target vertex $v$, it takes the
midpoint $m$ of the opposite edge and tests sources
$g=m+\varepsilon(v-m)/\lVert v-m\rVert$ for a systematically enumerated
positive margin.  It retains a witness only when $\lVert g\rVert<R$,
$\lVert v-g\rVert\leq1000$, and all 31 points have been checked against the
closed forward half-plane.  The resulting JSON records the source,
normal, source radius, target distance, and second-detector distance for all
31 targets.  The elementary geometry explains the separation: after the
small inward shift, the two edge endpoints have negative projection, while
other points in the positive side start at the lattice distance
$\sqrt7s/2$ from $m$ (the script verifies this from coordinates), so the
chosen target is the sole reachable detector.  This proves indispensability
of this particular 31-point lattice, not global optimality among arbitrary
placements.  The witness can be embedded in a legal 10--16-source world by
placing additional sources on distinct channels; no claim about a 31-point
global minimum is made.

Finally, an outward source on the arena boundary has angular half-width at
most $\beta=\arctan(q/R)$, hence at $q=1000$ at least
$\lceil\pi/\beta\rceil=7$ outer detectors are needed.  This is a global
outer-detector lower bound only; it does not imply that 31 total arbitrary
placements are optimal.  Boundary points and forward tests are closed, so
zero projection counts as reception, while the positive source shift makes
the selected witness strict.  The same certificate covers sources on the
boundary and at lattice vertices; for the latter the six nearest neighbours
are within $R+s$ and one has positive projection for every direction.
